\documentclass[../main.tex]{subfiles}
\begin{document}

\chapter{Advanced Caches}
\lessoninfo{
  video={Lesson 14, videos 1--40},
  textbook={H\&P \cite{hennessy2017} §2.3, §B.1, §B.3},
  keywords={pipelined cache, VIPT cache, aliasing, way prediction, NMRU, pseudo-LRU, conflict miss, prefetching, loop interchange, miss-under-miss support, MSHR, cache hierarchy, global miss rate, inclusion property},
}

Lesson 12 described the basic cache and measured it by the average memory
access time, and Lesson 13 added the address translation that every memory
access needs in a system with virtual memory.  This lesson collects the
techniques that make real caches fast.  They are organized by the term of the
average memory access time that they reduce: the hit time (pipelined caches,
the placement of the cache relative to the TLB, way prediction and cheap
replacement policies), the miss rate (larger blocks, prefetching and loop
interchange) and the miss penalty (servicing overlapping misses and, above
all, multilevel cache hierarchies).

\section{Improving cache performance}
\label{sec:l14-improving}

The quantity to be improved is the average memory access time (AMAT) of
Lesson~12:\source{Lesson 14, videos 1--2; H\&P §2.3.}
\begin{equation}
  \text{AMAT} = \text{hit time} + \text{miss rate} \times \text{miss penalty} .
  \label{eq:l14-amat}
\end{equation}
Each of the three terms offers its own line of attack, and
\cref{tab:l14-techniques} lists the techniques of this lesson by the term
they target.  Few of them are free.  A technique that lowers one term often
raises another---a smaller cache has a shorter hit time but a higher miss
rate---and it is worthwhile only if AMAT as a whole decreases.

\begin{table}[htbp]
  \centering\small
  \caption{Techniques of this lesson, grouped by the term of
    \cref{eq:l14-amat} they reduce.}
  \label{tab:l14-techniques}
  \begin{tabular}{@{}lll@{}}
    \toprule
    Term & Techniques & Sections \\
    \midrule
    Hit time     & pipelined caches, VIPT caches,               & \labelcref{sec:l14-pipelined}--\labelcref{sec:l14-replacement} \\
                 & way prediction, NMRU and PLRU replacement    & \\
    Miss rate    & larger blocks, prefetching, loop interchange & \labelcref{sec:l14-three-cs}--\labelcref{sec:l14-loop-interchange} \\
    Miss penalty & miss-under-miss support, cache hierarchies   & \labelcref{sec:l14-overlap}--\labelcref{sec:l14-inclusion} \\
    \bottomrule
  \end{tabular}
\end{table}

\section{Pipelined caches}
\label{sec:l14-pipelined}

The most direct way to lower the hit time is to keep the cache small and
simple, for example direct-mapped: fewer lines have to be decoded and fewer
tags compared.\source{Lesson 14, videos 3--4; H\&P §2.3.}  A small cache,
however, misses more often, and a direct-mapped cache suffers more conflicts,
so the time saved on hits is easily lost again on misses.  The techniques of
this and the next four sections reduce the hit time, or its cost, without
giving up the miss rate.

Lesson 12 noted that an L1 hit takes one to three clock
cycles.\caution{The L1 latencies vendors publish are load-to-use, not the
array access alone: Skylake quotes 4 cycles, and 5 with an indexed address.
Intel optimization manual.}  If a lookup
that needs several cycles occupied the cache until it completed, each memory
access would have to wait for the previous one: the hit time seen by an
access would be its own lookup time plus this wait, and consecutive loads and
stores would stall the processor.  A cache whose lookup fits into one cycle
has no such problem.  An L1 cache that takes two or three cycles is usually
pipelined.

\begin{definition}[Pipelined cache]
  A \term{pipelined cache}[パイプライン化キャッシュ] divides a lookup into
  stages separated by latches, so that a lookup takes several clock cycles but
  a new lookup can start in every cycle.
\end{definition}

For a set-associative cache the lookup divides naturally into three steps,
each of which can become a stage (\cref{fig:l14-pipelined}):
\begin{enumerate}
  \item the index selects the set, and the valid bits and tags of its lines
    are read;
  \item the tags are compared with the tag of the address, which decides
    between hit and miss and identifies the matching way, and the read of the
    data begins;
  \item the data read completes, the data of the matching way is selected,
    and the block offset picks the requested bytes out of the block, which are
    delivered.
\end{enumerate}
A miss is known at the end of the second stage.\margin{Reading every way's
data beside the tags is what makes step~2 fast.  Large L2 and L3 caches read
the tags first and then only the matching way: slower, much less energy.}
Pipelining does not reduce
the number of cycles one lookup takes.  It prevents that number from limiting
the rate of accesses: while one access compares its tags, the next one already
reads its set, exactly as instructions overlap in the pipeline of Lesson 3,
and the clock period needs to accommodate only the slowest stage instead of
the whole lookup.

\begin{figure}[htbp]
  \centering
  % Three consecutive accesses to a set-associative cache whose lookup is
% split into three pipeline stages: a new access starts in every cycle.
\begin{tikzpicture}[x=1.95cm,y=1.0cm, font=\footnotesize\sffamily,
  st/.style={draw, minimum width=1.85cm, minimum height=0.9cm, inner sep=1pt,
    align=center, font=\scriptsize\sffamily},
  s1/.style={st, fill=ln-box},
  s2/.style={st, fill=ln-accent!15},
  s3/.style={st, fill=ln-tip!20}]
  \foreach \c in {1,...,5} { \node at (\c,0.75) {\c}; }
  \node[anchor=east] at (0.45,0.75) {\iflangja{サイクル}{cycle}};
  \foreach \i/\n in {0/1,1/2,2/3} {
    \node[anchor=east] at (0.45,-\i) {\iflangja{アクセス}{access}~\n};
    \node[s1] at (\i+1,-\i) {\iflangja{セット選択\\V・タグ読出し}{select set,\\read V and tags}};
    \node[s2] at (\i+2,-\i) {\iflangja{タグ比較，\\データ読出し開始}{compare tags,\\start data read}};
    \node[s3] at (\i+3,-\i) {\iflangja{読出し完了，\\ウェイ選択，\\オフセットで抽出}{finish read,\\select way,\\offset $\to$ bytes}};
  }
\end{tikzpicture}

  \caption{Three consecutive accesses to a cache whose lookup is split into
    three stages.  Each access takes three cycles, but a new one starts in
    every cycle.}
  \label{fig:l14-pipelined}
\end{figure}

\begin{example}
  \label{ex:l14-pipelined}
  The lookup logic of an L1 cache has a delay of \SI{0.75}{\nano\second},
  and the other stages of the processor allow a clock period of
  \SI{0.3}{\nano\second}.  Without pipelining, either the clock period grows
  to \SI{0.75}{\nano\second} or more for the whole processor, or each access
  occupies the cache for three cycles and memory accesses can start only in
  every third cycle.  Split into three stages of about
  \SI{0.25}{\nano\second} each, which leaves room for the latch overhead, the
  lookup fits the \SI{0.3}{\nano\second} clock: the data arrives three cycles
  after the access starts, and an access can start in every
  cycle.\margin{Pipelining gives one access per cycle; ports and banks give
  more.  Skylake's L1 data cache serves two loads and one store per cycle.
  Intel optimization manual.}
\end{example}

\section{Address translation and the hit time}
\label{sec:l14-translation}

With virtual memory (Lesson 13) the processor issues virtual addresses,
whereas data in main memory is located by physical
addresses.\source{Lesson 14, videos 5--7; H\&P §B.3.}  The cache can be
looked up after the translation, before it, or partly in parallel with it
(\cref{fig:l14-translation}), and the choice affects the hit time.

\begin{figure}[htbp]
  \centering
  % Where address translation sits relative to the cache lookup:
% (a) physically accessed, (b) virtually accessed, (c) VIPT cache.
\begin{tikzpicture}[x=1cm,y=1cm, font=\footnotesize\sffamily, >={Stealth[length=1.8mm]},
  blk/.style={draw, fill=ln-box, minimum height=0.6cm, inner sep=2pt},
  cch/.style={draw, fill=ln-accent!15, minimum height=0.7cm, minimum width=2.2cm, inner sep=2pt},
  fld/.style={draw, fill=white, minimum height=0.5cm, inner sep=1pt, font=\scriptsize\sffamily},
  lab/.style={font=\scriptsize\sffamily, text=ln-muted, align=left},
  ttl/.style={font=\footnotesize\sffamily, align=center}]
  % ---------------- (a) physically accessed
  \node (va) at (1.6,0) {\iflangja{仮想アドレス}{virtual address}};
  \node[blk, minimum width=1.4cm] (tlb) at (1.6,-1.2) {TLB};
  \node[cch] (c) at (1.6,-2.7) {\iflangja{キャッシュ}{cache}};
  \node (d) at (1.6,-3.8) {\iflangja{データ}{data}};
  \draw[->] (va) -- (tlb);
  \draw[->] (tlb) -- node[lab, right] {\iflangja{物理\\アドレス}{physical\\address}} (c);
  \draw[->] (c) -- (d);
  \node[ttl] at (1.6,-4.55) {\iflangja{(a) 物理アドレスキャッシュ}{(a) physically accessed}};
  % ---------------- (b) virtually accessed
  \node (vb) at (5.2,0) {\iflangja{仮想アドレス}{virtual address}};
  \node[cch] (cb) at (4.9,-2.7) {\iflangja{キャッシュ}{cache}};
  \node (db) at (4.9,-3.8) {\iflangja{データ}{data}};
  \node[blk, dashed, fill=white, minimum width=1.2cm] (tb) at (6.3,-1.2) {TLB};
  \node[lab, anchor=north] at (6.3,-1.5) {\iflangja{変換は\\ミス時のみ}{translation\\only on a miss}};
  \draw[->] (4.9,-0.25) -- (cb);
  \draw[->, dashed] (6.3,-0.25) -- (tb);
  \draw[->] (cb) -- (db);
  \node[ttl] at (5.3,-4.55) {\iflangja{(b) 仮想アドレスキャッシュ}{(b) virtually accessed}};
  % ---------------- (c) VIPT
  \node[fld, minimum width=1.3cm] (vpn) at (9.05,0) {VPN};
  \node[fld, minimum width=1.3cm, anchor=west] (idx) at (vpn.east) {\iflangja{インデックス}{index}};
  \node[lab, anchor=south] at (9.7,0.25) {\iflangja{仮想アドレス}{virtual address}};
  \node[blk, minimum width=1.2cm] (tc) at (9.05,-1.2) {TLB};
  \node[cch, minimum width=2.6cm] (cc) at (9.7,-2.7) {\iflangja{キャッシュ}{cache}};
  \node (dc) at (9.7,-3.8) {\iflangja{データ}{data}};
  \draw[->] (vpn) -- (tc);
  \draw[->] (tc) -- node[lab, left] {\iflangja{フレーム\\番号}{frame\\number}} (tc.south |- cc.north);
  \draw[->] (idx.south) -- (idx.south |- cc.north);
  \draw[->] (cc) -- (dc);
  \node[ttl] at (9.7,-4.55) {(c) VIPT};
\end{tikzpicture}

  \caption{Three ways to combine the TLB with the cache lookup.  In~(b)
    translation is needed only on a miss, so the TLB is not in the path of a
    hit (it is still consulted for the permissions); in~(c) it translates the
    virtual page number while the index already selects the set.}
  \label{fig:l14-translation}
\end{figure}

\subsection{Physically accessed caches}

A cache that is looked up with the physical address is a
\term{physically accessed cache}[物理アドレスキャッシュ].  Every access first
translates its virtual address in the TLB, and only the resulting physical
address is used for the lookup (\cref{fig:l14-translation}a).  The hit time
therefore includes the translation:
\begin{equation}
  \text{hit time} = \text{TLB lookup time} + \text{cache lookup time} .
  \label{eq:l14-pa-hit}
\end{equation}
The TLB is small and fast, but its delay is paid on every access, hits
included, and it adds to exactly the quantity we are trying to
reduce.\margin{The TLB has to be tiny to keep up: Skylake has 64 first-level
data TLB entries for 4~KB pages, and 1536 second-level entries behind them.
Intel optimization manual.}

\subsection{Virtually accessed caches}

\begin{definition}[Virtually accessed cache]
  A \term{virtually accessed cache}[仮想アドレスキャッシュ] takes both the
  index and the tag from the virtual address.  The address is translated only
  on a miss, when the block has to be fetched from main memory.
\end{definition}

On a hit the TLB is not in the path, and the hit time is just the cache
lookup time (\cref{fig:l14-translation}b).  This advantage comes with
problems.
\begin{itemize}
  \item \textbf{Permissions.}  The TLB entry of a page also records whether
    the page may be read, written or executed, and the permission has to be
    checked on every access.  In practice the TLB is therefore consulted on
    hits as well.
  \item \textbf{Context switches.}  Each process has its own virtual address
    space, so the same virtual address denotes different data in different
    processes.  When the operating system switches the processor from one
    process to another (a context switch), the blocks in the cache are
    identified by virtual addresses of the old process and must not be found
    by the new one.  Either the cache is flushed by clearing all valid bits,
    or every tag is extended with an identifier of the address space (H\&P
    §B.3).  Context switches occur many times per second, and with flushing
    the process that runs next starts each time with an empty cache and a
    burst of misses.
  \item \textbf{Aliasing}, the subject of \cref{sec:l14-aliasing}.
\end{itemize}

\subsection{Virtually indexed, physically tagged caches}

A combination of the two approaches keeps the short hit time and avoids the
flushes.

\begin{definition}[VIPT cache]
  A \term{virtually indexed physically tagged cache}[仮想インデックス物理タグキャッシュ]%
  \index{VIPT|see{virtually indexed physically tagged cache}} (VIPT cache)
  takes the index from the virtual address and the tag from the physical
  address.
\end{definition}

The index bits of the virtual address select the set at once
(\cref{fig:l14-vipt}).  While the valid bits and tags of the set are being
read, the TLB translates the virtual page number into the physical frame
number, and the tags of the set, which are physical, are compared with the
frame number from the TLB.
\begin{itemize}
  \item The hit time is essentially the cache lookup time.  The TLB lookup
    runs in parallel with the selection of the set and only has to deliver
    the frame number by the time the tags are compared, which a small TLB
    usually manages.
  \item The cache need not be flushed on a context switch.  A tag names the
    physical memory its block came from, and that remains correct whichever
    process is running.
  \item The TLB is consulted on every access anyway, so the permissions are
    checked at no extra cost.
\end{itemize}

\begin{figure}[htbp]
  \centering
  % Lookup in a VIPT cache: 32-bit addresses, 4 KB pages, 64-byte blocks,
% 64 sets of 2 lines (8 KB).  The index comes from the page offset of the
% virtual address; the TLB translates the VPN in parallel, and the frame
% number is compared with the (physical) tags of the set; a line hits only if
% its valid bit is set.
\begin{tikzpicture}[x=1cm,y=1cm, font=\footnotesize\sffamily, >={Stealth[length=1.8mm]}]
  % --- virtual address fields
  %% The Japanese index label is wider than its box in the English layout,
  %% so in Japanese the VPN/index boundary (and the page-offset bracket that
  %% starts there) sits further left; the TLB arrow at x=1.3 and the index
  %% line at x=6.2 stay inside their boxes.
  \iflangja{%
    \draw[fill=ln-accent!15] (0.6,0) rectangle (4.0,0.55);
    \node at (2.3,0.275) {VPN\ (31\,\dots\,12)};
    \draw[fill=ln-box] (4.0,0) rectangle (7.0,0.55);
    \node at (5.5,0.275) {インデックス\ (11\,\dots\,6)};
  }{%
    \draw[fill=ln-accent!15] (0.6,0) rectangle (4.6,0.55);
    \node at (2.6,0.275) {VPN\ (31\,\dots\,12)};
    \draw[fill=ln-box] (4.6,0) rectangle (7.0,0.55);
    \node at (5.8,0.275) {index\ (11\,\dots\,6)};
  }
  \draw[fill=white] (7.0,0) rectangle (9.4,0.55);
  \node at (8.2,0.275) {\iflangja{オフセット}{offset}\ (5\,\dots\,0)};
  \node[anchor=south west, font=\sffamily\scriptsize, text=ln-muted, inner sep=1pt] at (0.6,0.6) {\iflangja{仮想アドレス}{virtual address}};
  % page offset bracket
  \iflangja{%
    \draw[ln-muted] (4.02,0.62) -- (4.02,0.78) -- (9.38,0.78) -- (9.38,0.62);
    \node[anchor=south, font=\sffamily\scriptsize, text=ln-muted] at (6.7,0.78)
      {ページオフセット (11\,\dots\,0)};
  }{%
    \draw[ln-muted] (4.62,0.62) -- (4.62,0.78) -- (9.38,0.78) -- (9.38,0.62);
    \node[anchor=south, font=\sffamily\scriptsize, text=ln-muted] at (7.0,0.78)
      {page offset (11\,\dots\,0)};
  }
  % --- TLB
  \node[draw, fill=ln-box, minimum width=1.3cm, minimum height=0.6cm] (tlb) at (1.3,-1.3) {TLB};
  \draw[->] (1.3,0) -- (tlb.north);
  % --- two ways, four drawn sets
  \foreach \x in {2.4,6.6} {
    \fill[ln-accent!15] (\x,-1.55) rectangle (\x+3.4,-2.0);
    \node[font=\sffamily\scriptsize, text=ln-muted, anchor=south] at (\x+0.2,-1.1) {V};
    \node[font=\sffamily\scriptsize, text=ln-muted, anchor=south] at (\x+1.1,-1.1) {\iflangja{タグ}{tag}};
    \node[font=\sffamily\scriptsize, text=ln-muted, anchor=south] at (\x+2.6,-1.1) {\iflangja{データ}{data}};
    \foreach \i in {0,...,3} {
      \draw (\x,-1.1-0.45*\i) rectangle (\x+0.4,-1.55-0.45*\i);
      \draw (\x+0.4,-1.1-0.45*\i) rectangle (\x+1.8,-1.55-0.45*\i);
      \draw (\x+1.8,-1.1-0.45*\i) rectangle (\x+3.4,-1.55-0.45*\i);
    }
  }
  \node[font=\sffamily\scriptsize, text=ln-muted, anchor=north] at (5.0,-2.95) {\iflangja{ウェイ 0}{way 0}};
  \node[font=\sffamily\scriptsize, text=ln-muted, anchor=north] at (9.2,-2.95) {\iflangja{ウェイ 1}{way 1}};
  \foreach \i/\s in {0/0,1/1,2/{\dots},3/63} {
    \node[font=\sffamily\scriptsize, text=ln-muted, anchor=west] at (10.05,-1.325-0.45*\i) {\s};
  }
  \node[font=\sffamily\scriptsize, text=ln-muted, anchor=south west] at (10.0,-1.1) {\iflangja{セット}{set}};
  % --- index selects a set in both ways
  \draw (6.2,0) -- (6.2,-1.775);
  \fill (6.2,-1.775) circle (1.2pt);
  \draw[->] (6.2,-1.775) -- (5.8,-1.775);
  \draw[->] (6.2,-1.775) -- (6.6,-1.775);
  % --- comparators against the frame number from the TLB
  \node[draw, circle, fill=white, inner sep=1pt, minimum size=0.5cm] (c0) at (3.5,-3.7) {$=$};
  \node[draw, circle, fill=white, inner sep=1pt, minimum size=0.5cm] (c1) at (7.7,-3.7) {$=$};
  \draw[->] (3.5,-2.9) -- (c0.north);
  \draw[->] (7.7,-2.9) -- (c1.north);
  % valid bits enable the comparators
  \draw[->] (2.6,-2.9) -- (c0.north west);
  \draw[->] (6.8,-2.9) -- (c1.north west);
  \draw (tlb.south) -- (1.3,-4.9) -- (7.7,-4.9);
  \draw[->] (7.7,-4.9) -- (c1.south);
  \fill (1.3,-3.7) circle (1.2pt);
  \draw[->] (1.3,-3.7) -- (c0.west);
  \node[anchor=east, align=right, font=\sffamily\scriptsize, text=ln-muted] at (1.2,-2.7)
    {\iflangja{物理\\フレーム\\番号}{physical\\frame\\number}};
  \node[draw, fill=ln-box, minimum width=0.7cm, minimum height=0.5cm, inner sep=1pt, font=\sffamily\scriptsize] (or) at (5.6,-3.7) {OR};
  \draw[->] (c0.east) -- (or.west);
  \draw[->] (c1.west) -- (or.east);
  \draw[->] (or.south) -- (5.6,-4.3) node[below] {\iflangja{ヒット}{hit}};
\end{tikzpicture}

  \caption{Lookup in a 2-way VIPT cache with 4~KB pages, 64-byte blocks and
    64 sets (8~KB).  The index comes from the page offset of the virtual
    address; the tags are physical frame numbers and are compared with the
    output of the TLB.  A line hits if its valid bit is set and its tag
    matches.}
  \label{fig:l14-vipt}
\end{figure}

\section{Aliasing}
\label{sec:l14-aliasing}

Virtual memory allows the same physical memory to appear at more than one
virtual address.\source{Lesson 14, videos 8--11; H\&P §B.3.}

\begin{definition}[Aliasing]
  The situation in which two or more virtual addresses, in the same or in
  different address spaces, are mapped to the same physical address is called
  \term{aliasing}[エイリアシング].  The virtual addresses are then aliases of
  each other.
\end{definition}

Aliasing is common: processes that share memory see the same frames, usually
at different virtual addresses, and a single process can have one frame
mapped at two addresses.  In a virtually accessed cache two aliases are
simply two different addresses.  They generally have different indexes and
tags, so the cache can hold two copies of the same physical block, one per
alias, and nothing connects the copies.

\begin{example}
  \label{ex:l14-alias}
  The virtual addresses V1 and V2 are aliases of the physical address P, and
  in a virtually accessed cache they have different indexes.  The memory at P
  contains the value~5.  \Cref{tab:l14-alias} follows four accesses.  The
  last one reads~5 through V2, although~7 has been written to the same
  memory through V1.  With the write-through policy main memory would
  hold~7, but the stale copy for V2 would still be returned.
\end{example}

\begin{table}[htbp]
  \centering\small
  \caption{Two aliases in a virtually accessed cache with the write-back
    policy (\cref{ex:l14-alias}): contents of the two cached copies and of
    main memory after each access.}
  \label{tab:l14-alias}
  \begin{tabular}{@{}clllll@{}}
    \toprule
    Step & Access          & Result          & Copy for V1 & Copy for V2 & Memory at P \\
    \midrule
    1    & read V1         & miss            & 5           & ---         & 5 \\
    2    & read V2         & miss            & 5           & 5           & 5 \\
    3    & write 7 to V1   & hit             & 7           & 5           & 5 \\
    4    & read V2         & hit, returns 5  & 7           & 5           & 5 \\
    \bottomrule
  \end{tabular}
\end{table}

To be correct, every write would have to find all other cached copies of the
same physical block and update or invalidate them, and the cache has no
efficient way of locating them.  Aliasing is thus a further problem of
virtually accessed caches.\sidenote{The operating system can prevent the
second copy instead: page coloring maps aliases only at virtual addresses
that agree in the index bits, so they select the same set.  The price is less
freedom in choosing which frame backs a page (H\&P §B.3).}

\subsection{Aliasing in VIPT caches}

In a VIPT cache the tag is physical, so two copies of a block would carry the
same tag.  They can nevertheless end up in different sets, because the index
is taken from the virtual address.  Whether that is possible depends on where
the index bits lie.  Translation replaces the virtual page number by the frame
number and leaves the page offset unchanged; aliases therefore always have the
same page offset, although their virtual page numbers differ.
\begin{itemize}
  \item If the index includes bits of the virtual page number, two aliases can
    select different sets, the cache can hold two copies, and the problem of
    \cref{ex:l14-alias} remains.
  \item If the index consists of page-offset bits only, aliases always select
    the same set.  There they find a single line with the physical tag of
    their block, so at most one copy exists and aliasing is harmless.  The
    index bits of the virtual address are then also those of the physical
    address, and the cache behaves exactly like a physically accessed cache,
    only faster.
\end{itemize}

\subsection{Avoiding aliasing}

A VIPT cache is thus free of aliasing when its block offset and index fit
into the page offset:
\begin{equation}
  \text{offset bits} + \text{index bits} \;\le\; \text{page offset bits} .
  \label{eq:l14-vipt-bits}
\end{equation}
The left side is $\log_2(\text{sets} \times \text{block size})$ and the right
side $\log_2(\text{page size})$, so the condition reads
$\text{sets} \times \text{block size} \le \text{page size}$.  Multiplying both
sides by the associativity~$N$ gives a bound on the cache size:
\begin{equation}
  \text{cache size} \;\le\; N \times \text{page size} .
  \label{eq:l14-vipt-size}
\end{equation}
The block size has cancelled out: only the associativity and the page size
limit a VIPT cache.\tip{Count bits, not bytes: with 4~KB pages only bits
0--11 survive translation, so offset bits $+$ index bits $\le 12$.}

\begin{example}
  \label{ex:l14-vipt-size}
  With 4~KB pages the page offset has 12 bits.  A VIPT cache with 64-byte
  blocks (6 offset bits) can use at most 6 index bits, i.e.\ 64 sets.  With 2
  ways this is the 8~KB cache of \cref{fig:l14-vipt}; with 4 ways it is
  16~KB.  A 2-way cache of 16~KB would need 128 sets, and its seventh index
  bit would be bit~12 of the virtual address, the lowest bit of the virtual
  page number, so aliases could be placed in different sets.
\end{example}

\subsection{Real VIPT caches}

The page size is fixed by the architecture and the operating system, and is
typically 4~KB.  By \cref{eq:l14-vipt-size}, an L1 cache that is to be both
VIPT and free of aliasing can therefore grow only by becoming more
associative: a direct-mapped cache is limited to 4~KB, a 32~KB cache needs 8
ways.  The Intel Core~2, Nehalem, Sandy Bridge and Haswell processors all
have 8-way L1 data caches of 32~KB, exactly at this bound.\margin{A larger
page raises the bound: Apple's M-series uses 16~KB pages, and the M1
performance core has a 128~KB L1 data cache.  Apple silicon CPU optimization
guide.}  Real L1 caches
are consequently more associative than their hit time alone would justify,
which makes the techniques of the next two sections important.

\begin{keypoint}
  A VIPT cache indexes with the virtual address while the TLB translates in
  parallel, and compares physical tags.  It has nearly the hit time of a
  virtually accessed cache, needs no flush on a context switch, and is immune
  to aliasing as long as $\text{cache size} \le N \times \text{page size}$.
\end{keypoint}

\section{Way prediction}
\label{sec:l14-way-prediction}

Associativity moves the hit time and the miss rate in opposite
directions.\source{Lesson 14, videos 12--15; H\&P §2.3.}  More ways per set
mean fewer conflicts and a lower miss rate---and for a VIPT cache a larger
permissible size (\cref{eq:l14-vipt-size})---but a lookup has to compare more
tags and select among more lines, so the hit time grows.  A direct-mapped
cache has the shortest hit time and the highest miss rate.  The goal is the
miss rate of a set-associative cache with the hit time of a direct-mapped one.

\begin{definition}[Way prediction]
  A set-associative cache with \term{way prediction}[ウェイ予測] guesses
  which line of the selected set holds the requested block and first checks
  only that line.  If its tag
  matches, the access completes in the hit time of a direct-mapped cache;
  otherwise the other lines of the set are checked as in an ordinary
  set-associative cache.
\end{definition}

The first check of an $N$-way cache of size~$C$ examines one line per set,
so it behaves like a direct-mapped cache of size $C/N$ with the same sets,
and the hit rate of such a cache estimates how often the guess is right.  If
the guess is the line accessed most recently in the set, the estimate is
exact: that line holds the block accessed most recently in the set, which is
precisely the block a direct-mapped cache of size $C/N$ would hold.  A wrong
guess costs the quick check plus the normal lookup, and a miss is detected
only after the normal lookup as well.  If $p$ is the fraction of accesses in
which the predicted line hits, $t_1$ the time to check one line and $t_N$ the
lookup time of the whole set, the average lookup time is
\begin{equation}
  t = p\, t_1 + (1-p)(t_1 + t_N) = t_1 + (1-p)\, t_N .
  \label{eq:l14-way-prediction}
\end{equation}

\begin{example}
  \label{ex:l14-way-prediction}
  An L1 cache is followed by an L2 cache that gives it a miss penalty of
  \SI{20}{\nano\second}.  Three L1 designs are compared.  The 4~KB
  direct-mapped cache is the subset of the 8-way cache that way prediction
  checks first, so the predicted line hits in $p = \SI{80}{\percent}$ of all
  accesses, and the remaining \SI{20}{\percent} (hits in other lines and all
  misses) take $0.8 + 1.5 = \SI{2.3}{\nano\second}$.
  \begin{center}\small
    \begin{tabular}{@{}lccl@{}}
      \toprule
      L1 cache & Lookup time & Miss rate & AMAT \\
      \midrule
      4~KB direct-mapped     & \SI{0.8}{\nano\second} & \SI{20}{\percent} & $0.8 + 0.20\times 20 = \SI{4.8}{\nano\second}$ \\
      32~KB 8-way            & \SI{1.5}{\nano\second} & \SI{5}{\percent}  & $1.5 + 0.05\times 20 = \SI{2.5}{\nano\second}$ \\
      32~KB 8-way, predicted & 0.8 or \SI{2.3}{\nano\second} & \SI{5}{\percent} & $1.1 + 0.05\times 20 = \SI{2.1}{\nano\second}$ \\
      \bottomrule
    \end{tabular}
  \end{center}
  By \cref{eq:l14-way-prediction} the average lookup time with prediction is
  $0.8 + 0.2\times 1.5 = \SI{1.1}{\nano\second}$.  Way prediction keeps the
  miss rate of the 8-way cache and brings its average lookup time close to
  that of the direct-mapped cache.
\end{example}

\section{Replacement policy and hit time}
\label{sec:l14-replacement}

The replacement policy affects the hit time as well as the miss
rate.\source{Lesson 14, videos 16--19; H\&P §B.1.}  Two of the policies of
Lesson 12 mark the extremes.
\begin{description}
  \item[Random] replacement does nothing on a hit and picks a victim at random
    on a miss.  It adds nothing to the hit time, but it sometimes evicts a
    block that is about to be used, so its miss rate is higher.
  \item[LRU] has a lower miss rate, but it keeps a $\log_2 N$-bit counter in
    every line, and an access, hit or miss, may change the counters of all
    $N$ lines of the set.  This activity on every access costs logic, energy
    and hit time, the more so the higher the associativity.
\end{description}
The policies below approach the miss rate of LRU with far less bookkeeping.

\subsection{Not most recently used}

\begin{definition}[NMRU]
  The \term{not most recently used}[NMRU]%
  \index{NMRU|see{not most recently used}} (NMRU) policy records for each set
  only which line was accessed most recently, and on a miss evicts one of the
  other lines, chosen at random.
\end{definition}

By temporal locality, the most recently used block is the one most likely to
be accessed next.  Evicting it is the most harmful choice random replacement
can make, and NMRU excludes exactly that choice.  The cost is one
$\log_2 N$-bit pointer per set, and an access changes at most this pointer: an
8-way cache needs 3 bits per set instead of the $8 \times 3 = 24$ bits of LRU.
NMRU misses somewhat more often than LRU, because its victim may be a block
that was used only a little less recently than the most recent one.

\subsection{Pseudo-LRU}

\begin{definition}[PLRU]
  The \term{pseudo-LRU}[疑似 LRU]\index{PLRU|see{pseudo-LRU}} (PLRU) policy
  keeps one bit per line, initially~0.\caution{Elsewhere \enquote{pseudo-LRU} usually names the tree variant:
  $N-1$ bits per set steer a binary tree to the victim.  The bit per line here
  is a different approximation.}
  An access sets the bit of its line;
  if this sets the last bit of the set that was still~0, the bits of all
  other lines are cleared.  A miss evicts a line whose bit is~0, chosen at
  random.
\end{definition}

The bit of a line tells whether the line has been accessed since the last
reset.  The bit of the most recently used line is always set, so the
candidates of PLRU are always among those of NMRU, and PLRU is at least as
good as NMRU.  How closely it approximates LRU depends on how many bits are
set.  Right after a reset only the most recently used line is excluded, as
in NMRU.  With $N-1$ bits set, the only candidate is the line whose last
access lies furthest in the past, which is exactly the choice of LRU.  The
cost is $N$ bits per set, and an access sets a single bit, except at a reset.

\begin{example}
  \label{ex:l14-replacement}
  A 4-way set holds the blocks A, B, C and D in lines 0--3; D was used most
  recently, and the PLRU bits were last reset by the access to~D.
  \Cref{tab:l14-replacement} follows the set under the three policies.  At
  the miss on~E, LRU must evict A, NMRU may evict A, C or D, and PLRU may
  evict A or C; here all three evict A.  If the next access, to F, also
  misses, LRU evicts C, NMRU may evict B, C or D, and PLRU, whose only
  line with bit~0 is line~2, also evicts C.
\end{example}

\begin{table}[htbp]
  \centering\small
  \caption{State kept by LRU, NMRU and PLRU for one 4-way set
    (\cref{ex:l14-replacement}): the LRU order from least to most recently
    used block, the most recently used line recorded by NMRU, and the PLRU
    bits of lines 0--3.}
  \label{tab:l14-replacement}
  \begin{tabular}{@{}llcccc@{}}
    \toprule
    Access & Result & Lines 0--3 & LRU order & NMRU & PLRU bits \\
    \midrule
    (initial) &      & A B C D & A B C D & 3 & 0 0 0 1 \\
    B         & hit  & A B C D & A C D B & 1 & 0 1 0 1 \\
    D         & hit  & A B C D & A C B D & 3 & 0 1 0 1 \\
    B         & hit  & A B C D & A C D B & 1 & 0 1 0 1 \\
    E         & miss & E B C D & C D B E & 0 & 1 1 0 1 \\
    \bottomrule
  \end{tabular}
\end{table}

\begin{keypoint}
  Work done on every access lengthens every hit.  Way prediction checks one
  line before the others, and NMRU and PLRU replace the counters of LRU by one
  pointer per set (NMRU) or one bit per line (PLRU); each keeps most of the
  miss-rate benefit of associativity or LRU at a fraction of the hit-time
  cost.
\end{keypoint}

\section{Causes of misses}
\label{sec:l14-three-cs}

Reducing the miss rate starts with knowing why misses
occur.\source{Lesson 14, video 20; H\&P §B.3; \cite{hill1989}.}  Misses are
classified into three kinds, often called the three Cs.

\begin{definition}[Compulsory, capacity and conflict misses]
  A \term{compulsory miss}[初期参照ミス] is a miss on the first access to a
  block; it occurs in any cache, even an infinitely large one.

  A \term{capacity miss}[容量性ミス] is a miss that would occur even in a
  fully associative cache of the same size: the cache cannot hold all the
  blocks the program is using, and the block has been evicted for lack of
  space.

  A \term{conflict miss}[競合性ミス] is a miss that occurs only because of
  limited associativity: a fully associative cache of the same size would
  still hold the block, but in the actual cache it was evicted by another
  block competing for the same set.\margin{A fourth C
  is added in multiprocessors: coherence misses, caused by another core's write
  invalidating the block (Lesson 19).}
\end{definition}

The three counts can be measured by simulating the same sequence of accesses
on three caches.  The misses of an infinite cache are the compulsory misses;
a fully associative cache of the actual size adds the capacity misses; the
actual cache adds the conflict misses.

\begin{example}
  \label{ex:l14-three-cs}
  For one program, an infinite cache misses \num{1500} times, a fully
  associative 32~KB cache \num{6000} times, and the actual 4-way 32~KB cache
  \num{7800} times.  Of the \num{7800} misses, \num{1500} are compulsory,
  $\num{6000} - \num{1500} = \num{4500}$ are capacity misses and
  $\num{7800} - \num{6000} = \num{1800}$ are conflict misses.
\end{example}

The classification points to the remedies.  A larger cache reduces capacity
misses, and a higher associativity or a better replacement policy reduces
conflict misses, but all of these lengthen the hit time.  The following
sections present techniques that reduce the miss rate without lengthening
hits.

\section{Larger cache blocks}
\label{sec:l14-block-size}

For a cache of a given size, a larger block size affects the miss rate in two
opposite ways.\source{Lesson 14, videos 21--22; H\&P §B.3.}
\begin{itemize}
  \item Where spatial locality is good, the data around an accessed address
    is used soon.  A larger block brings more of it into the cache with a
    single miss, and later accesses to that data, which would otherwise have
    missed, now hit.
  \item Where spatial locality is poor, most of a larger block is never used.
    The cache holds fewer blocks, hence fewer of the blocks that are actually
    needed, and capacity and conflict misses increase.
\end{itemize}
As the block size grows, the miss rate therefore first falls and then rises
again (\cref{fig:l14-block-size}).  The minimum lies at a larger block size
for a larger cache, for example at about 64~bytes for a 4~KB cache and at
about 256~bytes for a 256~KB cache: a large cache has room for many blocks
even when they are large, so the space wasted on unused data hurts it less.

\begin{figure}[htbp]
  \centering
  % Miss rate as a function of the block size for a 4 KB and a 256 KB cache
% (shape only, no data).  The best block sizes follow the lecture: about
% 64 bytes for the small and about 256 bytes for the large cache.
\begin{tikzpicture}[x=1cm,y=1cm, font=\footnotesize\sffamily, >={Stealth[length=1.8mm]}]
  \draw[->] (0,0) -- (7.6,0);
  \node[font=\footnotesize\sffamily] at (3.8,-0.75) {\iflangja{ブロックサイズ（バイト）}{block size (bytes)}};
  \draw[->] (0,0) -- (0,4.4) node[above, inner sep=2pt] {\iflangja{ミス率}{miss rate}};
  \foreach \x/\l in {0.8/16,2.0/32,3.2/64,4.4/128,5.6/256,6.8/512} {
    \draw (\x,0) -- (\x,-0.08) node[below, font=\scriptsize\sffamily] {\l};
  }
  \draw[thick, ln-accent] plot[smooth, tension=0.7] coordinates
    {(0.5,3.95) (0.8,3.6) (2.0,3.0) (3.2,2.82) (4.4,3.1) (5.6,3.65) (6.8,4.2)};
  \draw[thick, ln-tip] plot[smooth, tension=0.7] coordinates
    {(0.5,2.25) (0.8,1.95) (2.0,1.35) (3.2,1.0) (4.4,0.82) (5.6,0.74) (6.8,0.9) (7.2,1.02)};
  \fill[ln-accent] (3.2,2.82) circle (1.6pt);
  \fill[ln-tip] (5.6,0.74) circle (1.6pt);
  \node[anchor=west, text=ln-accent] at (7.0,4.2) {\iflangja{4~KB キャッシュ}{4~KB cache}};
  \node[anchor=west, text=ln-tip] at (7.3,1.02) {\iflangja{256~KB キャッシュ}{256~KB cache}};
  \node[font=\scriptsize\sffamily, text=ln-muted, align=center] at (1.9,2.3)
    {\iflangja{近傍へのアクセスが\\ヒットする}{accesses to neighbours\\now hit}};
  \node[font=\scriptsize\sffamily, text=ln-muted, align=center] at (5.3,2.3)
    {\iflangja{使われないデータが\\ラインを占める}{unused data\\occupies lines}};
  \draw[->, ln-muted, >={Stealth[length=1.4mm]}] (1.4,1.87) -- (2.9,1.87);
  \draw[->, ln-muted, >={Stealth[length=1.4mm]}] (4.6,1.87) -- (6.2,1.87);
\end{tikzpicture}

  \caption{Miss rate as a function of the block size for a 4~KB and a 256~KB
    cache (shape only).  The dots mark the best block sizes.}
  \label{fig:l14-block-size}
\end{figure}

\section{Prefetching}
\label{sec:l14-prefetching}

A miss can be avoided altogether if the block is brought into the cache
before the program asks for it.\source{Lesson 14, videos 23--26; H\&P
§2.3.}

\begin{definition}[Prefetching]
  Loading a block into the cache before it is accessed, because it is
  expected to be accessed soon, is called \term{prefetching}[プリフェッチ].
\end{definition}

Prefetching is a guess, and its effect depends on whether the guess is right.
\begin{itemize}
  \item A correct guess turns the first access to the block into a hit.  The
    fetch from memory still takes place, but earlier, while the processor is
    doing other work, so the miss penalty is not paid by the access.
  \item A wrong guess wastes the fetch.  Worse, the prefetched block has taken
    the place of another block, which may be needed later and then misses: the
    wrong guess \emph{pollutes} the cache.
\end{itemize}
Prefetching pays off only if most guesses are right.

\subsection{Prefetch instructions}

The program often knows which data it will access next.  A
\term{prefetch instruction}[プリフェッチ命令] is an instruction of the ISA
that asks the processor to bring the block containing a given address into
the cache; it changes no data.  The compiler or the programmer places it in
front of accesses that are known in advance, typically in loops over
arrays \cite{callahan1991}:
\begin{lstlisting}[language=C, numbers=none]
for (i = 0; i < n; i++) {
  prefetch(&a[i + D]);    /* D iterations ahead */
  sum = sum + a[i];
}
\end{lstlisting}
The difficulty is the prefetch distance~$D$.
\begin{itemize}
  \item If $D$ is too small, the prefetch has not completed when the loop
    reaches element $i+D$; the access still waits for the block, and little
    is gained.
  \item If $D$ is too large, the block arrives long before it is used.  It
    occupies a line in the meantime and may be evicted again before it is
    needed.
\end{itemize}
The right distance is about the time a fetch takes divided by the time of one
iteration.  Both depend on the machine, so a program compiled with a
distance that suits one processor prefetches too late or too early on
another.\caution{A prefetch distance is tuned for one machine.  A faster
processor with the same memory needs a larger distance.}

\begin{example}
  \label{ex:l14-prefetch-distance}
  A prefetched block is in the cache \SI{50}{\nano\second} after the prefetch
  instruction, and one iteration of the loop above takes
  \SI{4}{\nano\second}.  The distance must be at least $50/4 = 12.5$, so a
  distance of 16 leaves a small margin.  On a processor that executes each
  iteration in \SI{2}{\nano\second} with the same memory, the loop reaches
  element $i+16$ only \SI{32}{\nano\second} after its prefetch,
  \SI{18}{\nano\second} before the block arrives.  The loop then has to wait
  for memory: roughly every 16th access waits \SI{18}{\nano\second}, so 16
  iterations take about \SI{50}{\nano\second} instead of
  \SI{32}{\nano\second}, and the faster processor runs the loop at only about
  \SI{3.1}{\nano\second} per iteration.  It needs a distance of at least 25.
\end{example}

\subsection{Hardware prefetching}

In \term{hardware prefetching}[ハードウェアプリフェッチ] the processor itself
observes the sequence of memory accesses and guesses which blocks will be
needed.  The program need not be changed, and the guesses adapt to the
machine the program runs on.  Three designs exploit different access patterns
(\cref{fig:l14-prefetchers}).
\begin{description}
  \item[Stream buffer] A \term{stream buffer}[ストリームバッファ]
    \cite{jouppi1990} is a sequential prefetcher.  When blocks are accessed
    in order, it fetches the next few blocks into a small buffer beside the
    cache, from where an
    access takes them without a memory access, and it fetches further blocks
    as the buffer is used up.
  \item[Stride prefetcher] A \term{stride prefetcher}[ストライドプリフェッチャ]
    detects accesses that advance by a fixed distance, the \emph{stride}, as
    when a loop visits every $k$-th element of an array, and prefetches the
    address one or more strides ahead.
  \item[Correlating prefetcher] A
    \term{correlating prefetcher}[相関プリフェッチャ] records, for a block,
    which block was accessed after it, and when the block is accessed again
    prefetches the recorded successor.  It handles patterns without a regular
    spacing, such as a linked list whose nodes are scattered in memory and
    that is traversed repeatedly.
\end{description}

\begin{figure}[htbp]
  \centering
  % Access patterns exploited by three hardware prefetchers.  Numbers give
% the order of the accesses; P marks the block that is prefetched.
\begin{tikzpicture}[x=1cm,y=1cm, font=\footnotesize\sffamily, >={Stealth[length=1.5mm]}]
  \def\lnw{0.58}
  % \lnrow{y}{label}
  \def\lnrow#1#2{%
    \node[anchor=east, align=right] at (3.2,#1-0.29) {#2};
    \foreach \b in {0,...,11} {
      \draw[ln-rule, fill=white] ({3.4+\b*\lnw},#1) rectangle ++(\lnw,-0.58);
    }
  }
  % \lnacc{y}{block}{order}  accessed block
  \def\lnacc#1#2#3{%
    \draw[fill=ln-accent!25] ({3.4+#2*\lnw},#1) rectangle ++(\lnw,-0.58);
    \node at ({3.4+(#2+0.5)*\lnw},#1-0.29) {#3};
  }
  % \lnpre{y}{block}  prefetched block
  \def\lnpre#1#2{%
    \draw[fill=ln-tip!30] ({3.4+#2*\lnw},#1) rectangle ++(\lnw,-0.58);
    \node at ({3.4+(#2+0.5)*\lnw},#1-0.29) {P};
  }
  % block numbers
  \foreach \b in {0,...,11} {
    \node[font=\scriptsize\sffamily, text=ln-muted] at ({3.4+(\b+0.5)*\lnw},0.3) {\b};
  }
  \node[anchor=east, font=\scriptsize\sffamily, text=ln-muted] at (3.2,0.3) {\iflangja{ブロック}{block}};
  % (1) stream buffer: sequential
  \lnrow{0}{\iflangja{ストリームバッファ}{stream buffer}}
  \lnacc{0}{2}{1}\lnacc{0}{3}{2}\lnacc{0}{4}{3}
  \lnpre{0}{5}\lnpre{0}{6}
  \node[anchor=west, font=\scriptsize\sffamily, text=ln-muted] at (10.45,-0.29) {\iflangja{次のブロック}{next blocks}};
  % (2) stride prefetcher
  \lnrow{-1.0}{\iflangja{ストライド\\プリフェッチャ}{stride\\prefetcher}}
  \lnacc{-1.0}{1}{1}\lnacc{-1.0}{4}{2}\lnacc{-1.0}{7}{3}
  \lnpre{-1.0}{10}
  \node[anchor=west, font=\scriptsize\sffamily, text=ln-muted] at (10.45,-1.29) {\iflangja{ストライド 3}{stride 3}};
  % (3) correlating prefetcher
  \lnrow{-2.0}{\iflangja{相関\\プリフェッチャ}{correlating\\prefetcher}}
  \lnacc{-2.0}{8}{1}\lnacc{-2.0}{2}{2}\lnacc{-2.0}{5}{3}
  \draw[->, ln-accent] ({3.4+8.5*\lnw},-2.62) to[bend left=28] ({3.4+2.5*\lnw},-2.62);
  \draw[->, ln-accent] ({3.4+2.5*\lnw},-2.62) to[bend right=22] ({3.4+5.5*\lnw},-2.62);
  \node[anchor=west, font=\scriptsize\sffamily, text=ln-muted, align=left] at (10.45,-2.29)
    {\iflangja{8$\to$2，2$\to$5 を\\記録}{records\\8$\to$2, 2$\to$5}};
  % legend
  \draw[fill=ln-accent!25] (3.4,-3.55) rectangle ++(0.3,-0.3);
  \node[anchor=west, font=\scriptsize\sffamily] at (3.75,-3.7) {\iflangja{アクセス（数字は順序）}{accessed (number = order)}};
  \draw[fill=ln-tip!30] (7.4,-3.55) rectangle ++(0.3,-0.3);
  \node[anchor=west, font=\scriptsize\sffamily] at (7.75,-3.7) {\iflangja{プリフェッチ}{prefetched}};
\end{tikzpicture}

  \caption{Access patterns exploited by the three hardware prefetchers.
    After the three accesses shown, the stream buffer fetches the following
    blocks, the stride prefetcher the block one stride ahead, and the
    correlating prefetcher block~2 the next time block~8 is accessed.}
  \label{fig:l14-prefetchers}
\end{figure}

\section{Loop interchange}
\label{sec:l14-loop-interchange}

The compiler can also lower the miss rate by changing the order in which a
program accesses its data.\source{Lesson 14, video 27; H\&P §2.3.}  C stores
a two-dimensional array row by row: \texttt{m[j][i]} and
\texttt{m[j][i+1]} are adjacent, whereas \texttt{m[j+1][i]} lies a whole row
further on.  The following loop nest walks the array column by column:
\begin{lstlisting}[language=C, numbers=none]
double m[2000][1000];
for (i = 0; i < 1000; i++)
  for (j = 0; j < 2000; j++)
    m[j][i] = 2 * m[j][i];
\end{lstlisting}
Each access of the inner loop jumps \num{1000} elements, \num{8000} bytes,
beyond the previous one.  There is no spatial locality: every access of the
inner loop touches a different block, and the 2000 blocks visited by one pass
of the inner loop are far more than an L1 cache can hold, so nearly every
access misses.

\begin{definition}[Loop interchange]
  The compiler transformation that exchanges the order of nested loops, so
  that the innermost loop runs over the last index of the array and the data
  is accessed sequentially, is called \term{loop interchange}[ループ交換].
\end{definition}

After interchange the outer loop runs over \texttt{j} and the inner loop over
\texttt{i}.  The inner loop now visits adjacent elements, one block serves
several consecutive accesses, and a sequential or stride prefetcher can
anticipate the next block.  The compiler may interchange the loops only if it
can prove that the result is unchanged.  Here each iteration reads and writes
only its own element, so the iterations can be executed in any order.

\begin{keypoint}
  Prefetching can avoid even the miss on the first access to a block, but it
  helps only if its guesses are both correct and timely.  Loop interchange
  attacks misses from the program side, by making the order of accesses match
  the layout of the data in memory.
\end{keypoint}

\section{Overlapping misses}
\label{sec:l14-overlap}

The miss penalty can be larger than the time needed to fetch a
block.\source{Lesson 14, videos 28--30.}  An out-of-order processor (Lesson 7)
has many instructions in progress, and several of its loads and stores may
access memory within a few cycles.  Fetching a missing block takes many
cycles, and in the meantime another access may miss as well.

\begin{definition}[Overlap miss]
  An \term{overlap miss}[オーバーラップミス] is a cache miss that occurs
  while an earlier miss is still being serviced.
\end{definition}

After a load misses, an out-of-order processor continues with other
instructions.  The load cannot commit, however, so the processor eventually
runs out of ROB entries or reservation stations and stops; a second miss
issued before that point overlaps the first.  How much it costs depends on
the cache.
\begin{itemize}
  \item A \term{blocking cache}[ブロッキングキャッシュ] does nothing else
    while a miss is being serviced.  A later access is not even looked up
    until the earlier miss completes, so an overlapping miss pays the
    remaining time of the earlier miss in addition to its own penalty
    (\cref{fig:l14-overlap}a).
  \item A \term{non-blocking cache}[ノンブロッキングキャッシュ] serves hits
    while a miss is in progress (\emph{hit-under-miss}) and can also start
    further misses.  Several requests are then in progress in memory at the
    same time, which exploits \emph{memory-level parallelism}, provided that
    memory can service requests in parallel.
\end{itemize}

\begin{definition}[Miss-under-miss support]
  A cache with \term{miss-under-miss support}[ミスアンダーミス対応] can
  service several misses at the same time: it keeps track of each miss in
  progress and sends their requests to the next level of the memory hierarchy
  concurrently.
\end{definition}

Each miss in progress is recorded in a
\term{miss status handling register}[ミス状態管理レジスタ]%
\index{MSHR|see{miss status handling register}} (MSHR).  An access that
misses in the cache first compares its block address with the MSHRs.
\begin{description}
  \item[No match] The access is a new miss.  It allocates a free MSHR, which
    records the block and the instruction to be woken up, and the request is
    sent to memory.
  \item[Match] The block is already on its way; this is a
    \term{half-miss}[ハーフミス].  No second request is sent, and the
    instruction is only added to the matching MSHR.
\end{description}
When the data arrives, it is placed in the cache, all instructions recorded in
the MSHR are woken up, and the MSHR is released.  Already two MSHRs give a
large improvement over a blocking cache, four are better still, and because
memory accesses take so long, even 16 to 32 MSHRs continue to
help.\margin{Real machines track fewer: Intel cores from Nehalem to Skylake
have ten line-fill buffers, so at most ten L1 data misses are outstanding per
core.  Intel optimization manual.}

\begin{figure}[htbp]
  \centering
  % Two overlapping misses with a miss penalty of 20 cycles; the second miss
% occurs 6 cycles after the first.  (a) The cache services one miss at a
% time (blocking cache); (b) the cache has miss-under-miss support.  The red line marks the
% moment the second miss occurs.
\begin{tikzpicture}[x=0.215cm,y=1cm, font=\footnotesize\sffamily,
  svc/.style={draw, fill=ln-accent!25, minimum height=0.42cm, inner sep=0pt, font=\scriptsize\sffamily},
  lab/.style={anchor=east, font=\footnotesize\sffamily}]
  % \lnbar{x0}{x1}{y}{style}{text}
  \def\lnbar#1#2#3#4#5{%
    \draw[#4] (#1,#3+0.21) rectangle (#2,#3-0.21);
    \node[font=\scriptsize\sffamily] at ({(#1+#2)/2},#3) {#5};
  }
  % ---- (a) one miss at a time
  \node[anchor=west] at (-5,0.75) {\iflangja{(a) ブロッキングキャッシュ}{(a) blocking cache}};
  \node[lab] at (-1,0) {\iflangja{ミス 1}{miss 1}};
  \lnbar{0}{20}{0}{fill=ln-accent!25}{\iflangja{処理}{service}}
  \node[lab] at (-1,-0.6) {\iflangja{ミス 2}{miss 2}};
  \lnbar{6}{20}{-0.6}{fill=ln-box, draw=ln-rule}{\iflangja{待ち}{wait}}
  \lnbar{20}{40}{-0.6}{fill=ln-accent!25}{\iflangja{処理}{service}}
  \draw[ln-caution, thick] (6,-0.25) -- (6,-0.95);
  % ---- (b) miss-under-miss support
  \node[anchor=west] at (-5,-1.45) {\iflangja{(b) ミスアンダーミス対応}{(b) miss-under-miss support}};
  \node[lab] at (-1,-2.2) {\iflangja{ミス 1}{miss 1}};
  \lnbar{0}{20}{-2.2}{fill=ln-accent!25}{\iflangja{処理}{service}}
  \node[lab] at (-1,-2.8) {\iflangja{ミス 2}{miss 2}};
  \lnbar{6}{26}{-2.8}{fill=ln-accent!25}{\iflangja{処理}{service}}
  \draw[ln-caution, thick] (6,-2.45) -- (6,-3.15);
  % ---- time axis
  \draw[->] (0,-3.45) -- (42.5,-3.45);
  \foreach \t in {0,10,20,30,40} {
    \draw (\t,-3.45) -- (\t,-3.53) node[below, font=\scriptsize\sffamily] {\t};
  }
  \node[anchor=west, font=\scriptsize\sffamily] at (43,-3.45) {\iflangja{サイクル}{cycles}};
\end{tikzpicture}

  \caption{Two misses with a miss penalty of 20 cycles; the second occurs 6
    cycles after the first (vertical line).  (a)~In a blocking cache the
    second miss takes 34 cycles instead of 20.  (b)~With miss-under-miss
    support each miss takes 20 cycles.}
  \label{fig:l14-overlap}
\end{figure}

Miss-under-miss support costs additional hardware.  It pays off in
processors that issue many memory accesses concurrently, above all
out-of-order processors.

\section{Cache hierarchies}
\label{sec:l14-hierarchy}

The miss penalty of an L1 cache is dominated by the access time of main
memory.\source{Lesson 14, videos 31--34; H\&P §B.3.}  It can be reduced with
the same idea that led to the L1 cache in the first place: place another
cache between L1 and main memory.  A miss in L1 then looks up the
\term{L2 cache}[L2 キャッシュ], which is larger and slower than L1 but still
much faster than main memory, and only misses in L2 go further.

\begin{definition}[Cache hierarchy]
  A \term{cache hierarchy}[キャッシュ階層] consists of several levels of
  cache, L1, L2, L3, \ldots, each looked up only when the previous level
  misses (\cref{fig:l14-hierarchy}).  The level immediately before main
  memory is the \term{last-level cache}[最終レベルキャッシュ]%
  \index{LLC|see{last-level cache}} (LLC).
\end{definition}

\begin{figure}[htbp]
  \centering
  % A three-level cache hierarchy: each level is consulted only on a miss in
% the level before it; the miss penalty of a level is the average access
% time of everything behind it.
\begin{tikzpicture}[x=1cm,y=1cm, font=\footnotesize\sffamily, >={Stealth[length=1.8mm]},
  lvl/.style={draw, fill=ln-accent!15, align=center, inner sep=2pt},
  lab/.style={font=\scriptsize\sffamily, text=ln-muted}]
  \node[draw, fill=ln-box, minimum width=1.5cm, minimum height=0.9cm, align=center] (cpu) at (0,0)
    {\iflangja{プロセッサ}{processor}};
  \node[lvl, minimum width=1.0cm, minimum height=0.7cm] (l1) at (2.2,0) {L1};
  \node[lvl, minimum width=1.4cm, minimum height=1.0cm] (l2) at (4.3,0) {L2};
  \node[lvl, minimum width=1.9cm, minimum height=1.3cm] (l3) at (6.9,0) {L3\\(LLC)};
  \node[draw, fill=ln-box, minimum width=2.1cm, minimum height=1.6cm, align=center] (mm) at (10.0,0)
    {\iflangja{主記憶}{main\\memory}};
  \draw[->] (cpu) -- (l1);
  \draw[->] (l1) -- node[lab, above] {\iflangja{ミス}{miss}} (l2);
  \draw[->] (l2) -- node[lab, above] {\iflangja{ミス}{miss}} (l3);
  \draw[->] (l3) -- node[lab, above] {\iflangja{ミス}{miss}} (mm);
  % miss penalties
  \draw[ln-muted] (5.95,-0.9) -- (5.95,-1.05) -- (11.05,-1.05) -- (11.05,-0.9);
  \node[lab, anchor=north] at (8.5,-1.05) {\iflangja{L2 のミスペナルティ}{miss penalty of L2}};
  \draw[ln-muted] (3.6,-1.6) -- (3.6,-1.75) -- (11.05,-1.75) -- (11.05,-1.6);
  \node[lab, anchor=north] at (7.3,-1.75) {\iflangja{L1 のミスペナルティ}{miss penalty of L1}};
  % size / speed trend
  \draw[<->, ln-muted] (1.6,1.15) -- (11.0,1.15);
  \node[lab, anchor=south west] at (1.6,1.18) {\iflangja{小さく高速}{smaller, faster}};
  \node[lab, anchor=south east] at (11.0,1.18) {\iflangja{大きく低速}{larger, slower}};
\end{tikzpicture}

  \caption{A three-level cache hierarchy.  The miss penalty of each level is
    the average access time of the levels behind it.}
  \label{fig:l14-hierarchy}
\end{figure}

\subsection{AMAT of a cache hierarchy}

The miss penalty of a level is the average time it takes to obtain the block
from the levels behind it, which is the AMAT of the rest of the hierarchy.
Applying \cref{eq:l14-amat} level by level gives
\begin{equation}
  \begin{aligned}
    \text{AMAT} &= h_1 + m_1 \times P_1 , \\
    P_1 &= h_2 + m_2 \times P_2 , \\
    &\;\;\vdots \\
    P_{\text{LLC}} &= \text{access time of main memory},
  \end{aligned}
  \label{eq:l14-amat-hierarchy}
\end{equation}
where $h_i$ is the hit time, $m_i$ the miss rate and $P_i$ the miss penalty
of level~$i$.  The miss rate $m_i$ is counted over the accesses that reach
level~$i$ (\cref{sec:l14-local-global}).

\subsection{L1 versus L2}

In \cref{eq:l14-amat-hierarchy} the hit time $h_1$ of L1 is paid by every
access, whereas the hit time $h_2$ of L2 is paid only by the fraction $m_1$
of accesses that miss in L1.  L1 is therefore designed for a short hit time,
with the techniques of the previous sections, while L2 can trade a longer
hit time for a lower miss rate.\margin{Skylake, per Intel's optimization
manual: L1 data 32~KB, 4 cycles; L2 256~KB, 12 cycles; L3 about 2~MB per
core, 44 cycles; main memory 60--100~ns, i.e.\ hundreds of cycles.}

\begin{example}
  \label{ex:l14-hierarchy}
  Main memory takes \SI{50}{\nano\second} per access.  A 32~KB cache has a
  hit time of \SI{1}{\nano\second} and a miss rate of \SI{10}{\percent}; a
  256~KB cache has a hit time of \SI{5}{\nano\second} and a miss rate of
  \SI{3}{\percent}.  Assume that in a hierarchy the 256~KB cache misses the
  same accesses as it would alone, i.e.\ $0.03/0.10 = \SI{30}{\percent}$ of
  the accesses that reach it.
  \begin{center}\small
    \begin{tabular}{@{}ll@{}}
      \toprule
      Configuration & AMAT \\
      \midrule
      No cache                      & \SI{50}{\nano\second} \\
      32~KB cache only              & $1 + 0.10\times 50 = \SI{6}{\nano\second}$ \\
      256~KB cache only             & $5 + 0.03\times 50 = \SI{6.5}{\nano\second}$ \\
      L1 32~KB, L2 256~KB           & $1 + 0.10\times(5 + 0.30\times 50) = \SI{3}{\nano\second}$ \\
      \bottomrule
    \end{tabular}
  \end{center}
  The hierarchy is twice as fast as either cache alone.  It combines the
  short hit time of the small cache, paid by every access, with the low miss
  rate of the large one, whose longer hit time is paid by only
  \SI{10}{\percent} of the accesses.
\end{example}

\section{Local and global miss rates}
\label{sec:l14-local-global}

In \cref{ex:l14-hierarchy} the L2 cache misses \SI{30}{\percent} of the
accesses that reach it, three times the rate of L1.\source{Lesson 14, videos
35--37; H\&P §B.3.}  This does not make L2 a poor cache.  L2 sees only the
accesses that L1 has missed, which are the ones with the weakest locality; the
accesses that are easy to serve never get past L1.  The miss rate of a lower
level must therefore be stated with its reference.

\begin{definition}[Local and global rates]
  The \term{local hit rate}[局所ヒット率] of a cache is the number of its hits
  divided by the number of accesses that reach it, and its
  \term{local miss rate}[局所ミス率] is $1 - \text{local hit rate}$.

  The \term{global miss rate}[大域ミス率] of a cache is the number of its
  misses divided by the number of all memory accesses of the processor, and
  its \term{global hit rate}[大域ヒット率] is $1 - \text{global miss rate}$.
\end{definition}

For L1 the local and global rates coincide.  For a lower level the accesses
that reach it are the misses of the level above, so the global miss rate of
level~$n$ is the product of the local miss rates:
\begin{equation}
  \text{global miss rate}_n = \prod_{i=1}^{n} \text{local miss rate}_i .
  \label{eq:l14-global-miss}
\end{equation}
The miss rates in \cref{eq:l14-amat-hierarchy} are local.  The global miss
rate of a level says what fraction of all accesses has to go beyond it, which
makes it the right measure for comparing levels and hierarchies.

\begin{definition}[Misses per 1000 instructions]
  The metric \term{misses per 1000 instructions}[1000 命令当たりミス数]%
  \index{MPKI|see{misses per 1000 instructions}} (MPKI) counts the misses of
  a cache per 1000 executed instructions.
\end{definition}

Like the global miss rate, MPKI counts all misses of a level, but relative to
the instructions instead of the memory accesses:
\begin{equation}
  \text{MPKI} = 1000 \times \text{global miss rate} \times
  \frac{\text{memory accesses}}{\text{instruction}} .
  \label{eq:l14-mpki}
\end{equation}
It thus states directly how often per instruction a program pays the miss
penalty of that level.\tip{About a third of instructions access memory in
typical code, so MPKI $\approx 330 \times$ global miss rate (H\&P instruction
mixes).}

\begin{example}
  \label{ex:l14-local-global}
  A program executes \num{1000000} instructions with \num{1500000} memory
  accesses.  L1 misses \num{90000} times and L2 \num{27000} times.
  \begin{itemize}
    \item L1: local and global miss rate
      $\num{90000}/\num{1500000} = \SI{6}{\percent}$; MPKI
      $\num{90000}/1000 = 90$.
    \item L2: local miss rate $\num{27000}/\num{90000} = \SI{30}{\percent}$
      (local hit rate \SI{70}{\percent}); global miss rate
      $\num{27000}/\num{1500000} = \SI{1.8}{\percent} = 0.06\times 0.30$
      (global hit rate \SI{98.2}{\percent}); MPKI $27$.
  \end{itemize}
\end{example}

\begin{keypoint}
  AMAT is computed with local miss rates.  Levels and hierarchies are compared
  with global miss rates or MPKI, because a low local hit rate is normal for a
  level that sees only the misses of the level above.
\end{keypoint}

\section{The inclusion property}
\label{sec:l14-inclusion}

A block in L1 may or may not also be present in L2.\source{Lesson 14, videos
38--40; H\&P §B.3; \cite{baer1988}.}  Three relationships between the
levels are possible:
\begin{description}
  \item[Inclusion] every block in L1 is also in L2.
  \item[Exclusion] no block in L1 is also in L2.
  \item[Neither] some blocks in L1 are also in L2 and others are not.
\end{description}

\begin{definition}[Inclusion property]
  A cache hierarchy has the \term{inclusion property}[包含性] if every block
  in a level is also present in the next, larger level.
\end{definition}

It might seem that a larger L2 with LRU replacement keeps every block of L1
automatically.  It does not.  Hits in L1 never reach L2, so a block that is
used often, and therefore stays in L1, looks unused to the LRU counters of L2;
eventually it becomes the least recently used block of L2 and is evicted
there, while it remains in L1.

\begin{example}
  \label{ex:l14-inclusion}
  L1 holds 2 blocks and L2 holds 3 blocks; both are fully associative with LRU
  replacement and initially empty, and a block that misses in both levels is
  loaded into both.  \Cref{tab:l14-inclusion} follows the accesses P, Q, P, R,
  P, S.  The later accesses to P hit in L1, so L2 last sees P at the first
  access.  At the miss on S, P is the least recently used block of L2 and is
  evicted, although it is still in L1.
\end{example}

\begin{table}[htbp]
  \centering\small
  \caption{A two-level hierarchy with LRU replacement in both levels
    (\cref{ex:l14-inclusion}).  Contents are listed from least to most
    recently used; the last access breaks inclusion.}
  \label{tab:l14-inclusion}
  \begin{tabular}{@{}lllll@{}}
    \toprule
    Access & L1   & L2   & L1 after & L2 after \\
    \midrule
    P      & miss & miss & P        & P \\
    Q      & miss & miss & P Q      & P Q \\
    P      & hit  & ---  & Q P      & P Q \\
    R      & miss & miss & P R      & P Q R \\
    P      & hit  & ---  & R P      & P Q R \\
    S      & miss & miss & P S      & Q R S \\
    \bottomrule
  \end{tabular}
\end{table}

\needspace{12\baselineskip}
Inclusion therefore has to be enforced if it is wanted.  One way is an
\emph{inclusion bit} in every line of L2, set while the block is also in L1;
L2 never chooses a line with this bit set as its victim.\sidenote{Both
choices are in use.  Intel's client L3 up to Skylake is inclusive and doubles
as a filter: a block that is not in L3 is in no core's cache, so the other
cores need not be asked (Lesson 19).  The L3 of Skylake-SP and of AMD's Zen
is non-inclusive and holds mainly the blocks evicted from the private caches,
so its capacity is not spent on copies.}

This lesson completes the discussion of caches in a single processor
(Lesson 19 returns to caches in multiprocessors).  Lesson 15 turns to the
memory technology from which caches and main memory are built.

\needspace{10\baselineskip}
\begin{keypoint}[Lesson summary]
  AMAT $=$ hit time $+$ miss rate $\times$ miss penalty, and each term has its
  own techniques.  The hit time is reduced by pipelining the lookup, by VIPT
  caches, which index with the virtual address while the TLB translates and
  avoid aliasing if $\text{cache size} \le N \times \text{page size}$, by way
  prediction, and by the cheap replacement policies NMRU and PLRU.  Misses
  are compulsory, capacity or conflict misses; the miss rate is reduced by a
  suitable block size, by prefetching (prefetch instructions, stream buffers,
  stride and correlating prefetchers) and by loop interchange.  The miss
  penalty is reduced by miss-under-miss support, which tracks the misses in
  progress in MSHRs, and by cache hierarchies, whose AMAT is computed level
  by level with local miss rates, while global miss rates and MPKI describe
  how many accesses go past a level.  With LRU
  in each level, inclusion does not hold automatically.
\end{keypoint}

\end{document}
